\section{Reproducible Computational Study}\label{sec:experiments}
\subsection{Data partitions and frozen decisions}
The primary study contains 24 development problems and 48 new test problems, generated with disjoint seed ranges beginning at 510000 and 810000. A separate training bank for a neural policy comparison begins at 610000. A problem is an entire noisy DOSY matrix, not an individual spectral pixel. The support configuration was frozen before test evaluation: maximum order four, screening parameter $0.01$, validation stride four, one paired standard-deviation allowance, and at most 350 nonlinear function evaluations per support refinement.

An initial development run preferred support size before order. A weak-species development case revealed fragmentation, so the frozen version prefers order first. Both development runs are preserved. There was no test-set sweep of the screening parameter. The primary results reaggregate saved outputs and redraw figures without refitting the selected models or tuning them to the illustrated ground truth. The additional post-review experiments in Appendices~\ref{app:revisionvalidation} and \ref{app:completevalidation} are newly executed and reported separately.

Each new problem has 32 acquisition coordinates and 1024 spectral channels uniformly spanning 10 to 0 ppm. In the primary atomic benchmark, $D_0=10^{-9}\,\mathrm{m^2\,s^{-1}}$ and fitted rates are restricted to $[0.1,15]D_0$. The native rate constraints in the external comparisons are specified in Appendix~\ref{app:externalcomparison}. The acquisition design is
\begin{equation}\label{eq:design}
 b_i=10^9 B\left(\frac{i}{31}\right)^\eta\ \mathrm{s\,m^{-2}},
 \quad i=0,\ldots,31,
\end{equation}
with independent $B\sim U(1.3,2.7)$ and $\eta\sim U(1.5,2.4)$. The component diffusion coefficients are positive and jointly shared across all spectral channels. Their amplitudes have correlated Gaussian spectral peaks, including an overlapping region near 7.73 ppm and other peaks with different mixture proportions. The exact generator is included with the evidence package.

Eight families contain six problems each: one, two, or three separated rates; a three-rate weak-component family; a close-rate family; proportional component spectra; noise-only controls; and broad distributions. The six problems per family use SNR 50, 100, and 200 twice each, with independently generated parameters and noise. For positive signals, the maximum unattenuated spectral intensity is normalized to one in the generator and $\sigma=1/\mathrm{SNR}$. Solvers receive observations in these units, not true component normalizations.

Nominal separated rates are $(0.65,2.3,7.1)D_0$, using the first one or two as appropriate, with independent log-rate perturbations of standard deviation 0.10. Close rates use $(0.8,1.05,5.5)D_0$ with log-rate perturbations of standard deviation 0.02. In the weak family the third spectrum is multiplied by 0.04 before overall normalization; this is a spectral-amplitude multiplier, not a guarantee of exactly 4\% total mixture mass. In the proportional family $A_{j\cdot}$ equals $1,0.7,0.45$ times the same spectrum. Broad controls replace each atom by a lognormal distribution with log-width 0.18, using 24-point Gauss--Hermite forward quadrature. The forward generator uses the lognormal on its full positive domain; its histogram reference is truncated and renormalized on the fitted interval for plotting and shape evaluation. These controls deliberately violate the fitted atomic model.

Every fourth acquisition, using zero-based indices $i\bmod4=3$, is reserved for external testing. The remaining 24 rows are available to the reconstruction. A signal mask thresholds the mean of the first six available low-$b$ rows at $4\sigma/\sqrt6$, followed by two one-dimensional dilation steps. Noise-only controls instead inspect the fixed window $4<\mathrm{ppm}<4.65$. All methods in a comparison receive the same mask. Conditional on that mask, external test noise remains independent. The reported errors assess retained spectral channels, not undetected parts of a full spectrum; losses caused by masking are not chemical absence certificates.

After the 24 reconstruction rows are sorted by increasing acquisition coordinate, the 3rd, 7th, 11th, 15th, 19th and 23rd rows form the internal validation set. The other 18 rows are used for candidate fitting. These are one-based positions within the reconstruction subset, not positions in the original 32-row acquisition. Groups are split where adjacent selected chemical-shift spacings exceed 1.6 times their 25th percentile. This rule assumes a substantially regular original spectral grid. It is a segmentation convention, not a learned chemical grouping. The final selected model uses all 24 reconstruction rows.

\subsection{Evaluation, denominators, and uncertainty summaries}
The primary quantitative panel is the 36 signal-present atomic problems. Six null cases are reported separately because normalized shape and true spectral leakage have no physical denominator there. Six broad cases are misspecification controls. An additional 24-problem archived bank contains finite-width diffusion profiles with one to three generating components, including close and proportional scenarios. It is transfer evidence and is never pooled with the new atomic bank.

Rates are paired to ground truth by a minimum-cost one-to-one assignment in absolute log distance; pairs beyond $0.25$ are then rejected. ``Order and rates recovered'' requires the correct total order and an accepted pair for every true rate. The tolerance corresponds to an approximate multiplicative factor of $e^{0.25}$, and is not a chemical identification standard. On finite-width controls, matching an atomic rate to a distribution center does not imply recovery of the distribution or a literal atom.

For each accepted pair $(j,k)$, define an evaluator-only meaningful spectral set
$\mathcal S_k=\{\ell:A^{\mathrm{true}}_{k\ell}>0.001\max_h A^{\mathrm{true}}_{kh}\}$ within the supplied mask. The leakage fraction is
\begin{equation}\label{eq:leakage}
 \ell_{\mathrm{spec}}=
 \frac{\sum_{(j,k)\ \mathrm{paired}}\sum_{\ell\notin\mathcal S_k}\widehat A_{j\ell}}
 {\sum_{(j,k)\ \mathrm{paired}}\sum_\ell\widehat A_{j\ell}}.
\end{equation}
The implementation assigns zero when its denominator is zero; therefore unmatched estimated mass is always reported alongside it, as a fraction of all estimated mass. Missing true rates are exposed by the recovery count. The $0.1\%$ threshold is an evaluation convention and never a fitting cutoff. It is also distinct from the display threshold in the DOSY figures.

External clean prediction error is $\|\widehat Y_U-Y_U^0\|_F/(\sigma\sqrt{|U|p})$ on the eight unused acquisitions $U$. It measures estimation error without adding a fresh noise variance; observed noisy-test errors are retained in the machine-readable results. Shape loss is $W_1$ in natural-log diffusion, evaluated after exporting both measures to the common 256-cell mass grid, normalizing per positive spectral channel, and weighting spectral channels by true mass only in the evaluator. A missing positive estimate receives the declared span penalty $\log150$; it is not described as an actual transport distance. Total estimated mass is saved separately.

We report problem-level averages with equal weight per problem. The primary table excludes the null cases; Appendix~\ref{app:legacy} retains the earlier 42-case convention, whose zero contributions dilute shape averages. Bootstrap intervals resample whole paired problems 10,000 times with seed 9925. They are descriptive pointwise intervals for this small heterogeneous simulation bank, not simultaneous tests or calibrated population coverage. Different spectral channels are not treated as independent experimental replicates.

\subsection{Neural policy comparison and numerical verification}
RAI-Net is the action-ranking policy of Section~\ref{sec:rainetmath}, acting on the adaptive distributional objective in Equation~\eqref{eq:raihaar}. It does not directly predict the number of species. Retraining used 80 DOSY problems split by whole problem into 64 training and 16 validation problems, producing 2921 recorded action examples. The policy has two hidden layers of 48 units with hyperbolic-tangent activations. Training lasted 400 epochs and used a transformed noise-standardized internal validation prediction-gain target. The selected checkpoint was epoch 400, so convergence of learning is not claimed. The classical objective and numerical acceptance checks remained in force.

The original and retrained policies were evaluated on the predefined first 24 new test problems and all 24 archived transfer problems. The new neural panel contains three problems per scenario: 18 signal-present atomic, three null, and three broad-distribution cases. These are smaller panels than the atomic study. The network results are an auxiliary comparison, not an additional order-recovery experiment or evidence for the atomic support rule.

Numerical verification covers NNLS agreement with an independent implementation, finite-difference checks of the variable-projection Jacobian, coherent weak-signal screening, acquisition partition boundaries, and amplitude-scaling/spectral channel-permutation behavior. A direct comparison with MATLAB R2026a Update 2 \texttt{lsqnonneg} used 100 amplitude columns; a full MATLAB interface call was also checked. Numerical logs and hashes accompany the frozen outputs. Pseudoinverse-based local rate precision, boundary hits, and a residual-based atomic-model warning are exported as diagnostics; none is interpreted as a selection-adjusted confidence interval or goodness-of-fit test.

All figures were generated by deterministic scientific plotting of saved arrays.
Dated source archives identify the earlier benchmark computations. The programming
and language assistance used in preparing this work is disclosed at the end of
the manuscript.

\subsection{Revision studies and experimental sampling units}\label{sec:extendeddesign}
Table~\ref{tab:revisionunits} identifies the independent inputs and repeated
conditions in each additional study. Methods are compared on the same observations;
additional policies, candidate orders and warm-start budgets do not increase the
number of independent samples. The supplement contains the configurations, seeds,
checkpoints and failed outcomes. In the convergence study, we fix the factor count
to examine optimization alone. The weak-component and experimental studies retain
automatic atomic order selection.

\begin{table}[htbp]\centering\small
\caption{Sampling units of the revision studies. Repeated conditions and
methods do not multiply the number of independent inputs.}\label{tab:revisionunits}
\begin{tabularx}{\linewidth}{@{}lX@{}}\toprule
Panel & Independent input and repeated conditions\\\midrule
Null & 300 noise matrices; 100 per noise scale; paired methods.\\
Weak & 96 truth/noise matrices; four amplitude factors, three SNRs, eight per cell.\\
Factorial & 36 existing atomic problems; paired $2\times2$ stages and production reference.\\
TRAIn-MF & Nine rank--SNR inputs; four warm budgets and two common starts.\\
Neural & 120 new test problems; five trained seeds and one classical control on each.\\
External & 47 frozen matrices after one adapter pilot; two stresses on 12 underlying truths.\\
Experimental & Four resin specimen series with three diffusion-time blocks each.\\
Public mixture & One 30-acquisition archived experiment; five paired estimators.\\\bottomrule
\end{tabularx}\end{table}

We studied four industrial-resin series: TP194-2623086, TP194-2623088,
TP194-2623084 and TP194-2619082. We retain the signed absorption spectra and fit
every channel within the positive support detected from the fitting acquisitions.
Each 23-acquisition block reserves five acquisitions for prediction; neither
these nor the internal validation rows enter support detection. All methods use
the same encoding and support. Coherent negative structure outside the support is
reported and left unestimated. Strong correlation in the processed background
means that its estimated noise scale cannot be interpreted as an iid Gaussian
calibration. We retain the HDO-derived encoding calibration, but its unquantified
systematic uncertainty prevents assigning a chemical confidence interval.
Appendix~\ref{app:experimentalrevision} describes the samples and the limits of
the available composition records.